% SPDX-License-Identifier: Apache-2.0
% covers: RemoteLink
\datasheet{RemoteLink}{The wire under a remote peripheral: each transaction one Ethernet frame, and each frame that answers it an answer.}

\dsfacts{\code{lib/parts/src/remote/eth.rs}}{\code{txhdl\_parts::remote::eth::RemoteLink}}%
{\code{//docs:flagship}}

\subsection*{Function}

\code{RemoteLink} sits between a \code{Remote}'s two channels and a
MAC's. An \code{Ask} taken on \code{ask} becomes one frame of 27 bytes
on \code{tx}; a frame arriving on \code{rx} that answers it becomes an
\code{Answer} on \code{back}. The MAC of \code{//docs:eth} is what
puts those bytes on a wire and takes them off, and neither end of the
peripheral knows the wire is there.

Nothing else on the wire is looked at, and nothing else has to be. A
frame is this peripheral's when it carries this protocol's type, the
answer kind, and this peripheral's device number, and any byte that
says otherwise ends the matter for that frame.

\subsection*{Parameters}

\begin{center}\footnotesize
\begin{tabular}{@{}lp{0.7\textwidth}@{}}
\toprule
Parameter & Meaning \\
\midrule
\code{DEV} & The device number. It goes in every frame this sends and
must match every frame it accepts, so several of these can share one
wire and one program, each answered by the number it was asked under.
It is a parameter rather than a register because a design knows how
many peripherals it has. The tables use 3, as \code{ex\_remote} does;
the board uses 1. \\
\bottomrule
\end{tabular}
\end{center}

\subsection*{The frame}

27 bytes, before whatever padding the MAC adds to reach the minimum
frame. An answer is read at the same offsets, so a program answers by
changing four bytes of what it was sent.

\begin{center}\footnotesize
\begin{tabular}{@{}llp{0.52\textwidth}@{}}
\toprule
Bytes & Field & What \\
\midrule
0 to 6 & destination & Broadcast. The peripheral does not know where
the program is, and the program may move. \\
6 to 12 & source & \code{02:00:00:00:00:\textit{dev}}, locally
administered. \\
12 to 14 & type & \code{0x88b5}, which IEEE leaves for local
experimental use, so no registered protocol can be mistaken for it. \\
14 & kind & \code{0x01} asks, \code{0x81} answers. \\
15 & device & \code{DEV}. \\
16 & tag & The transaction's, which the answer repeats. \\
17 & flags & Bit 0 is a write on an ask, an error on an answer. \\
18 to 22 & address & Most significant byte first. \\
22 to 26 & word & The word written, or the word answered. \\
26 & strobe & Which lanes a write means. \\
\bottomrule
\end{tabular}
\end{center}

\dstables{RemoteLink}

\subsection*{Behaviour}

\begin{itemize}
\item One frame goes out at a time. An \code{Ask} is taken when
nothing is going out, and its fields are held in registers while the
frame is spelled from them a byte per cycle that \code{tx} has room,
\code{at} saying which byte is next.
\item The bytes are a function of \code{at} and those registers
rather than a buffer, so the frame costs no memory.
\item A frame coming in is counted by \code{got}, and four bytes are
checked as they pass: the two of the type, the kind, and the device.
\item What the check sets is \code{alien} rather than a match, because
a register starts at zero and a frame is this peripheral's until a
byte of it says otherwise. It is cleared on the byte after the last,
so the next frame starts innocent.
\item The tag, the error bit and the four bytes of the word are
latched as they arrive, at bytes 16, 17 and 22 to 25.
\item An \code{Answer} is sent on \code{back} on the last byte of a
frame that nothing objected to. A frame that is not this
peripheral's ends with nothing sent, and its bytes are read and
dropped rather than left to block the MAC.
\item Sending and receiving are independent: a frame can arrive while
one is going out.
\end{itemize}

\subsection*{Verification}

\begin{itemize}
\item Five tests in \code{lib/parts/src/remote/eth.rs}, in
\code{//lib/parts:parts\_test}.
\code{a\_transaction\_leaves\_as\_a\_frame} asserts all 27 bytes of a
write, one by one.
\code{a\_read\_leaves\_with\_no\_word\_and\_no\_strobe} is the other
kind.
\code{an\_answer\_frame\_becomes\_an\_answer} is the way back.
\code{a\_frame\_that\_is\_not\_this\_devices\_is\_ignored} is another
device's frame and another protocol's.
\code{the\_bus\_is\_served\_by\_a\_program\_that\_speaks\_frames} puts
a \code{Remote} in front of it and serves the bus from frames alone.
\item \code{//lib/examples:ex\_remote} runs the whole chain, two MACs
and a wire between them, and the build simulates this netlist against
that run under nvc and Verilator:
\code{//docs:remote\_sim\_remote\_link\_remote\_link\_tb\_test}
and \code{//docs:remote\_vsim\_remote\_link\_test}.
\item \code{//tools/remote:remote\_test} asserts the same 27 bytes
from the Go end, against the same constants rather than against this
implementation, so the two spellings of the protocol are held apart.
\end{itemize}

\subsection*{Used by}

\code{ex\_remote}, and \code{Board}, where a \code{RemoteLink<1>}
carries the remote peripheral at \code{0x3300} out to the flagship's
Ethernet port.

\subsection*{Limits}

One frame out at a time and one transaction in flight, since that is
all a \code{Remote} asks for. The frame is fixed at 27 bytes and
carries one word, so a burst is as many frames as it has beats. There
is no retry and no sequence number beyond the tag: a frame the wire
loses is a transaction that goes unanswered, and what covers it is
the peripheral's timeout. Nothing authenticates a frame, so anything
on the wire that speaks this protocol and knows the device number can
write the peripheral's address space.
